\documentclass[../main.tex]{subfiles}

\begin{document}

    We are given that $F_{main}(x,y) = 3x^{3} + 2y^{3} + 2x^{2} + 2xy + 2y^{2} - 6x - 6y$, and the
    vector gradient $\nabla F_{main}$ and the stationary points are defined as \\
    \begin{center*}
        \begin{align*}
            \nabla F_{main} =
            \begin{pmatrix}
                9x^{2} + 4x + 2y - 6 \\
                6y^{2} + 2x + 4y - 6
            \end{pmatrix}
            =
            \begin{pmatrix}
                0 \\
                0
            \end{pmatrix}
        \end{align*}
    \end{center*}
    In order to find the stationary point of $F_{main}$, we need to solve $\nabla F_{main}(x,y) = 0$.
    
    \section{Deriving our iterative function}\label{sec:2d-nr-deriving-iterative-function}
    \begin{center*}
        \begin{align*}
            & \text{Let} \ \mathbf{x}_{n} =
            \begin{pmatrix}
            x_{n} \\
            y_{n}
            \end{pmatrix}, \
            \nabla F_{main}(\mathbf{x}_{n}) =
            \begin{pmatrix}
                \frac{\partial F}{\partial x} \\
                \frac{\partial F}{\partial y}
            \end{pmatrix} =
            \begin{pmatrix}
                f(x_{n}, y_{n}) \\
                g(x_{n}, y_{n})
            \end{pmatrix} \\
            & \text{then,} \ H(x_{n}, y_{n}) =
            \begin{pmatrix}
                \frac{\partial^{2} F}{\partial x^{2}} & \frac{\partial^{2} F}{\partial x \partial y} \\
                \frac{\partial^{2} F}{\partial y \partial x} & \frac{\partial^{2} F}{\partial y^{2}}
            \end{pmatrix} =
            \begin{pmatrix}
                18x + 4 & 2 \\
                2 & 12y + 4
            \end{pmatrix}
        \end{align*}
    \end{center*}
    We can rearrange the Newton-Raphson identity to equal $0$,
    \begin{center*}
        \begin{align*}
            & x_{n + 1} = x_{n} - \frac{f(x_{n})}{f'(x_{n})} \\
            & \implies x_{n + 1} \cdot f'(x_{n}) = x_{n} \cdot f'(x_{n}) - f(x_{n}) \\
            & \implies f'(x_{n})(x_{n + 1} - x_{n}) = -f(x_{n}) \\
            & \implies f(x_{n}) + f'(x_{n})(x_{n + 1} - x_{n}) = 0
        \end{align*}
    \end{center*}
    and then, with $\nabla F_{main}(x_{n}, y_{n}) = \mathbf{F}(\mathbf{x}_{n})$, we can work back to the
    Newton-Raphson identity, but with matrices this time.
    \begin{center*}
        \begin{align*}
            & \mathbf{F}(\mathbf{x}_{n}) +
            H(\mathbf{x}_{n})(\mathbf{x}_{n + 1} - \mathbf{x}_{n}) = 0 \\
            & \implies H(\mathbf{x}_{n})(\mathbf{x}_{n + 1} - \mathbf{x}_{n}) = -\mathbf{F}(\mathbf{x}_{n}) \\
            & \implies \mathbf{x}_{n + 1} - \mathbf{x}_{n} = -H^{-1}(\mathbf{x}_{n})\mathbf{F}(\mathbf{x}_{n}) \\
            & \implies \mathbf{x}_{n + 1} = \mathbf{x}_{n} - H^{-1}(\mathbf{x}_{n})\mathbf{F}(\mathbf{x}_{n})
        \end{align*}
    \end{center*}
    
    \subsection*{Finding $H^{-1}(\mathbf{x}_{n})$}\label{subsec:2d-nr-finding-h^-1}
    We know that for a given matrix
    \begin{center*}
        \begin{align*}
            & A = \begin{pmatrix}
                a & b \\
                c & d
            \end{pmatrix} \\
            & A^{-1} = \frac{1}{ad-bc}\begin{pmatrix}
                d & -b \\
                -c & a
            \end{pmatrix}
        \end{align*}
    \end{center*}
    So,
    \begin{center*}
        \begin{align*}
            & H(\mathbf{x}_{n}) =
            \begin{pmatrix}
                18x + 4 & 2 \\
                2 & 12y + 4
            \end{pmatrix} \\
            & H^{-1}(\mathbf{x}_{n}) =
            \frac{1}{(18x + 4)(12y + 4) - 2 \cdot 2}\begin{pmatrix}
                12y + 4 & -2 \\
                -2 & 18x + 4
            \end{pmatrix} \\
            & \implies H^{-1}(\mathbf{x}_{n}) = \frac{1}{216xy + 72x + 48y + 12}\begin{pmatrix}
                12y + 4 & -2 \\
                -2 & 18x + 4
            \end{pmatrix}
        \end{align*}
    \end{center*}
    
    \section{Using our iterative function}\label{sec:2d-nr-iterative-function}
    We now have that
    \begin{center*}
        \begin{align*}
            \mathbf{x}_{n + 1} = \mathbf{x}_{n} - H^{-1}(\mathbf{x}_{n})\mathbf{F}(\mathbf{x}_{n}) \\
        \end{align*}
    \end{center*}
    which is equivilent to
    \begin{center*}
        \begin{align*}
            \begin{pmatrix}
                x_{n + 1} \\
                y_{n + 1}
            \end{pmatrix} = \begin{pmatrix}
                x_{n} \\
                y_{n}
            \end{pmatrix} - \frac{1}{216xy + 72x + 48y + 12}\begin{pmatrix}
                12y + 4 & -2 \\
                -2 & 18x + 4
            \end{pmatrix}\begin{pmatrix}
                9x^{2} + 4x + 2y - 6 \\
                6y^{2} + 2x + 4y - 6
            \end{pmatrix}
        \end{align*}
    \end{center*}
    
    \subsection*{Solving $\nabla F_{main}(x,y) = 0$}\label{subsec:2d-nr-solving-f}
    The initial condition is
    \begin{center*}
        \begin{align*}
            & \mathbf{x}_{0} = \begin{pmatrix}
                -5 \\ -1
            \end{pmatrix}
        \end{align*}
    \end{center*}
    and iterating 10 times yields:
    \begin{center*}
        \begin{align*}
            & \mathbf{x}_{1} = \begin{pmatrix}
                -5 \\ -1
            \end{pmatrix} - \frac{1}{684}\begin{pmatrix}
                -8 & -2 \\
                -2 & -86
            \end{pmatrix}\begin{pmatrix}
                197 \\
                -14
            \end{pmatrix} = \begin{pmatrix}
                -2.7368 \\ -2.1842
            \end{pmatrix} \\
            & \mathbf{x}_{2} = \begin{pmatrix}
                -1.6975 \\ -1.7118
            \end{pmatrix} \\
            & \vdots \\
            & \mathbf{x}_{10} = \begin{pmatrix}
                -1.2542 \\ -1.5699
            \end{pmatrix}
        \end{align*}
    \end{center*}

    \subsection*{Verifying the root}
    To verify we have found the root, all we need to do is plug $\mathbf{x}_{10}$ into $\nabla F_{main}$:
    \begin{center*}
        \begin{align*}
            \nabla F_{main}(\mathbf{x}_{10}) = \begin{pmatrix}
                9(-1.2542)^{2} + 4(-1.2542) + 2(-1.5699) - 6 \\
                6(-1.5699)^{2} + 2(-1.2542) + 4(-1.5699) - 6
            \end{pmatrix} = \begin{pmatrix}
                \num{5.5876e-4} \\
                \num{-4.8394e-4}
            \end{pmatrix} \approx \begin{pmatrix}
                0 \\ 0
            \end{pmatrix}
        \end{align*}
    \end{center*}
    
    \subsection*{Nature of the stationary point}\label{subsec:2d-nr-stationary-point-nature}
    First, we have to find the value of $H(\mathbf{x}_{10})$,
    \begin{center*}
        \begin{align*}
            H(\mathbf{x}_{10}) = \begin{pmatrix}
                18(-1.2542) + 4 & 2 \\
                2 & 12(-1.5699) + 4
            \end{pmatrix} = \begin{pmatrix}
                -18.5756 & 2 \\
                2 & -14.8388
            \end{pmatrix}
        \end{align*}
    \end{center*}
    Then, we calculate the eigenvalues of this matrix as
    \begin{center*}
        \begin{align*}
            & \lambda_{1} = -19.4438, \\
            & \lambda_{2} = -13.9705
        \end{align*}
    \end{center*}
    Since they are both negative, the stationary point is a \textbf{local maximum}.

    \subsection*{Code used to find the root}\label{subsec:2d-nr-code}
    I ran the following in my terminal (usage covered in
    \hyperref[subsubsec:2d-nr-appendix]{2D Newton-Raphson} under
    \hyperref[subsec:usage-of-modules]{Usage of modules}):
    \begin{lstlisting}[style=console,label={lst:2d-nr}]
C:\User\CaseStudy> python main.py 2d-nr
Root: [-1.25417304 -1.56992898]
Root verification: [0. 0.]
Hessian eigenvalues: [-19.44380116 -13.9704613 ]
Stationary point nature: local maximum
    \end{lstlisting}

\end{document}
